\documentclass[../../../include/open-logic-section]{subfiles}

\begin{document}

\olfileid{sth}{ord-arithmetic}{add}
\olsection{د ترتيبي عددونو جمع}

فرض کړئ چې $\alpha$ او $\beta$ جمع کول غواړو۔ د $\beta$ يوه
کاپي د $\alpha$ د يوې کاپۍ سملاسي وروسته کېښودلے شو۔ (کاپيو ته اړ يو،
ځکه له \olref[ordinals][basic]{ordinalsaresubsets} څخه پوهېږو چې يا
$\alpha \subseteq \beta$ وي يا $\beta \subseteq \alpha$۔) په بداهتي
ډول دا داسې ده لکه لومړے د $\alpha$ په اوږدوالي ګامونه واخلو،
\emph{او بيا} د $\beta$ په اوږدوالي ګامونه واخلو۔ پايلې ته رسېدلې
لړۍ به ښه ترتيب شوې وي؛ نو د
\olref[ordinals][ordtype]{thmOrdinalRepresentation} له مخې له يوه
(يوازيني) ترتيبي عدد سره هم‌شکله ده۔ هماغه ترتيبي عدد د $\alpha$
او $\beta$ \emph{مجموعه} ګڼلے شو۔

دا د ترتيبي جمع تر شا بداهتي تصور دے۔ د کره تعريف دپاره لومړے
د سټونو د \emph{کاپيو} جوړولو تصور اخلو۔ دلته د $0$ او $1$ خپل‌سري
ټاګونه کاروو، څو څرګنده شي چې کوم څيز له کومه راغلے دے:

\begin{defn}\ollabel{defdissum}
د $A$ او $B$ \emph{ناګډه مجموعه} دا ده:
$A \disjointsum B = (A\times \{0\}) \cup (B \times \{1\})$۔
\end{defn}

اوس د ترتيبي عددونو پر جوړو يو ترتيب تعريفوو:

\begin{defn}
د هر ډول ترتيبي عددونو $\alpha_1, \alpha_2, \beta_1, \beta_2$
دپاره وايو چې:
\begin{align*}
	\tuple{\alpha_1, \alpha_2} \rlexless \tuple{\beta_1, \beta_2}\text{ هله او يوازې هله چې }&
	\text{يا $\alpha_2 \in \beta_2$ وي}\\
	& \text{يا هم $\alpha_2 = \beta_2$ او $\alpha_1 \in \beta_1$ دواړه وي}
\end{align*}
\end{defn}

دا \emph{برعکس قاموسي} ترتيب دے، ځکه لومړے د جوړې دويم غړے او
بيا لومړے غړے پرتله کېږي۔ اوس په ياد ولرئ چې $\alpha \ordplus \beta$
مو د $\alpha$ د کاپۍ ورپسې د $\beta$ د کاپۍ د ترتيب ډول غوښت۔ د دې
دپاره وايو:

\begin{defn}\ollabel{defordplus}
د هر ډول ترتيبي عددونو $\alpha$، $\beta$ دپاره د دوی مجموعه دا ده:
$\alpha \ordplus \beta = \ordtype{\alpha \disjointsum \beta, \rlexless}$۔
\end{defn}
\noindent
پام وکړئ چې دلته مو د نښو په کارونه کښې لږ تساهل کړے دے؛ په کره
ډول بايد «$\Setabs{\tuple{x,y}\in \alpha\disjointsum\beta}{x \rlexless y}$»
د «$\rlexless$» پر ځاے وليکو۔ خو د لنډيز دپاره به لاندې همدا
تساهل روان ساتو۔

لاندې پايله، له \olref[ordinals][ordtype]{thmOrdinalRepresentation}
سره يوځاے، ښيي چې زموږ تعريف په سمه توګه جوړ شوے دے:

\begin{lem}\ollabel{ordsumlessiswo}
$\tuple{\alpha \disjointsum \beta, \rlexless}$ ښه ترتيب دے،
د هر ډول ترتيبي عددونو $\alpha$ او $\beta$ دپاره۔
\end{lem}

\begin{proof}
په څرګنده $\rlexless$ پر $\alpha \disjointsum \beta$ تړلې اړيکه ده۔
د دې د ښه‌بنسټه والي د ښودلو دپاره ناتش
$X \subseteq \alpha \disjointsum \beta$ وټاکئ۔ $Y$ دې د $X$
هغه فرعي سټ وي چې د غړو دويم مختص ئې تر ټولو کوچنے وي، يعنې
$Y = \Setabs{\tuple{\gamma, i} \in X}{(\forall \tuple{\delta, j} \in X)i \leq j}$۔
اوس د $Y$ هغه غړے وټاکئ چې لومړے مختص ئې تر ټولو کوچنے وي۔
\end{proof}
\noindent
په دې توګه د ترتيبي جمع يو ښه ښکاره تعريف لرو۔ يوه د تعجب وړ
نه ده پايله دا ده ($1 = \{0\}$ په ياد ولرئ، د
\olref[z][infinity-again]{defnomega} له مخې):
\begin{prop}
$\alpha \ordplus  1 = \ordsucc{\alpha}$، د هر ترتيبي عدد
$\alpha$ دپاره۔
\end{prop}

\begin{proof}
د $f$ هم‌شکلتيا له $\ordsucc{\alpha} = \alpha \cup \{\alpha\}$ څخه تر
$\alpha\disjointsum1 = (\alpha \times \{0\})
\cup (\{0\} \times \{1\})$ پورې وګورئ؛ دا
$f(\gamma) = \tuple{\gamma, 0}$ د $\gamma \in \alpha$ دپاره،
او $f(\alpha) = \tuple{0, 1}$ ټاکي۔
OLSTH-013: د سرچينې په ښودل شوي قيمت کښې دواړه برخې له مخکې
ټاګ شوې دي؛ د دويمې ناګډې جمع پر ځاے د همدغو برخو عادي اتحاد
ښودل شوے دے۔ د f قيمتونه نه دي بدل شوي۔
\end{proof}
\noindent
سربېره پر دې، په اسانه ښودل کېږي چې جمع ځينې بازګشتي شرطونه
پوره کوي:

\begin{lem}\ollabel{ordadditionrecursion}
د هر ډول ترتيبي عددونو $\alpha, \beta$ دپاره:
\begin{align*}
	\alpha\ordplus 0 &= \alpha\\
	\alpha \ordplus  (\beta\ordplus 1) &= (\alpha \ordplus  \beta) \ordplus  1\\
	\alpha  \ordplus  \beta &= \supstrict_{\delta < \beta}(\alpha \ordplus  \delta) && \text{که $\beta $ حدي ترتيبي عدد وي}
\end{align*}
\end{lem}

\begin{proof}
هر حالت جلا ګورو؛ لومړے:
\begin{align*}
	\alpha \ordplus  0
	& = \ordtype{(\alpha \times \{0\}) \cup (0 \times \{1\}), \rlexless} \\
	&= \ordtype{(\alpha \times \{0\}) \cup \emptyset, \rlexless}\\
	&= \alpha\\
	\alpha \ordplus (\beta \ordplus  1)
	%&= \ordtype{(\alpha\times \{0\}) \cup (\ordtype{\beta \ordplus  1}\times \{1\}), \rlexless} \\
	&= \ordtype{(\alpha\times \{0\}) \cup (\ordsucc{\beta}\times \{1\}), \rlexless} \\
	&= \ordtype{(\alpha\times \{0\}) \cup (\beta \times \{1\}), \rlexless} \ordplus 1\\
	&= (\alpha \ordplus  \beta) \ordplus  1
\end{align*}
OLSTH-014: د سرچينې په دويمه معادله کښې د تش سټ او د يوغړي
سټ نښې سره ګډې شوې وې؛ د صفر او يوه د حاصل ضرب پر ځاے تش سټ
راغلے دے۔ لومړۍ او وروستۍ معادلې عين دي۔
اوس دې $\beta \neq \emptyset$ حدي وي۔ که $\delta < \beta$ وي،
نو $\delta\ordplus 1 < \beta$ هم دے؛ له دې امله
$\alpha \ordplus  \delta$ د $\alpha \ordplus  \beta$ يوه
مناسبه پيلنۍ برخه ده۔ نو $\alpha \ordplus  \beta$ د
$X = \Setabs{\alpha \ordplus  \delta}{\delta < \beta}$ کلکه
پاسنۍ کرانه ده۔ سربېره پر دې، که
$\alpha \leq \gamma < \alpha \ordplus  \beta$ وي، نو په ښکاره
$\gamma = \alpha \ordplus \delta$ د کوم $\delta < \beta$
دپاره دے۔ نو
$\alpha \ordplus  \beta = \supstrict_{\delta< \beta}(\alpha\ordplus \delta)$۔
\end{proof}

خو دلته يوه په زړه پورې خبره ده۔ د ترتيبي جمع د تعريف دپاره مو
\emph{پر ځاے} يوازې د نامتناهيت هاخوا بازګښت قضيه کارولے شوه،
او د \olref{ordadditionrecursion} هماغه بازګشتي معادلې مو ايښودلے
شوې (البته «$\ordsucc{\beta}$» په کارولو، د
«$\beta \ordplus 1$» پر ځاے)۔

نو د ترتيبي عددونو عمليې په دوو بېلو لارو تعريفولے شو۔ په
\emph{ترکيبي} لاره يو ښه ترتيب شوے سټ په ښکاره جوړوو او د هغه
د ترتيب ډول اخلو۔ يا ئې \emph{بازګشتي} لاره تعريفولے شو، يوازې
د بازګښت معادلې په ايښودلو۔ که کار سم شوے وي، پايله يوه ده؛ ځکه
\olref[ordinals][ordtype]{thmOrdinalRepresentation} تضمينوي چې
دغه بازګشتي معادلې \emph{يوازيني} ترتيبي عددونه ټاکي۔

په ډېرو اړخونو کښې د ترتيبي عددونو حساب د طبيعي عددونو د جمع
په شان دے۔ د بېلګې په توګه، لاندې ثابتولے شو:

\begin{lem}\ollabel{ordinaladditionisnice}
که $\alpha, \beta, \gamma$ ترتيبي عددونه وي، نو:
\begin{enumerate}
	\item\ollabel{ordaddition1} که $\beta < \gamma$ وي، نو $\alpha
	\ordplus  \beta < \alpha \ordplus  \gamma$
	\item\ollabel{ordaddition2} که $\alpha \ordplus  \beta =
	\alpha\ordplus \gamma$ وي، نو $\beta = \gamma$
	\item\ollabel{ordaddition3} $\alpha \ordplus  (\beta \ordplus
	\gamma) = (\alpha \ordplus  \beta) \ordplus  \gamma$، يعنې
	جمع اتحادي ده
	\item\ollabel{ordaddition4} که $\alpha \leq \beta$ وي، نو $\alpha
	\ordplus  \gamma \leq \beta \ordplus \gamma$
\end{enumerate}
\end{lem}

\begin{proof}
\olref{ordaddition3} ثابتوو او پاتې برخې د تمرين دپاره پرېږدو۔
ثبوت د $\gamma$ په اړه د ساده نامتناهيت هاخوا استقرا له لارې،
د \olref{ordadditionrecursion} په کارولو کېږي۔ که $\gamma = 0$ وي:
\[
(\alpha \ordplus  \beta) \ordplus  0 = \alpha \ordplus  \beta  = \alpha \ordplus  (\beta \ordplus  0)
\]
که $\gamma = \delta\ordplus 1$ وي، د استقرا له مخې فرض کړئ چې
$(\alpha \ordplus  \beta) \ordplus  \delta = \alpha \ordplus  (\beta \ordplus
\delta)$؛ اوس \olref{ordadditionrecursion} درې وارې کاروو:
\begin{align*}
	(\alpha \ordplus  \beta) \ordplus  (\delta \ordplus  1) & = ((\alpha \ordplus  \beta) \ordplus  \delta)\ordplus 1\\
	& = (\alpha \ordplus  (\beta \ordplus  \delta)) \ordplus  1\\
	& = \alpha \ordplus  ((\beta \ordplus  \delta)\ordplus 1)\\
	& = \alpha \ordplus  (\beta \ordplus  (\delta\ordplus 1))
\end{align*}
که $\gamma$ حدي ترتيبي عدد وي، د استقرا له مخې فرض کړئ چې که
$\delta \in \gamma$ وي، نو
$(\alpha \ordplus  \beta) \ordplus  \delta =
\alpha \ordplus  (\beta \ordplus  \delta)$؛ اوس:
\begin{align*}
	(\alpha \ordplus  \beta) \ordplus  \gamma & = \supstrict_{\delta < \gamma}((\alpha \ordplus  \beta) \ordplus  \delta) \\
	&= \supstrict_{\delta < \gamma}(\alpha \ordplus  (\beta \ordplus  \delta))\\
	&= \alpha \ordplus  \supstrict_{\delta < \gamma}(\beta \ordplus  \delta)\\
	& = \alpha \ordplus  (\beta \ordplus  \gamma)
\end{align*}
\end{proof}

\begin{prob}
د \olref[sth][ord-arithmetic][add]{ordinaladditionisnice} پاتې
برخې ثابتې کړئ۔
\end{prob}

تر دې ځايه ترتيبي جمع بايد ډېره اشنا ښکاره شي۔ خو يوه مهمه
لاره شته چې ترتيبي جمع پکې د طبيعي عددونو له جمع سره \emph{نه}
ورته کېږي۔

\begin{prop}\ollabel{ordsumnotcommute}
ترتيبي جمع تبادلي \emph{نه} ده؛ $1 \ordplus  \omega = \omega <
\omega \ordplus  1$۔
\end{prop}

\begin{proof}
پام وکړئ چې $1 \ordplus  \omega = \supstrict_{n < \omega} (1 \ordplus n)
= \omega \in \omega \cup \{\omega\} = \ordsucc{\omega} = \omega
\ordplus  1$۔
\end{proof}
\noindent
که څه هم دا په لومړي نظر حيرانولے شي، \emph{بايد مو حيران نه کړي}۔
له يوې خوا، په $1 \ordplus  \omega$ کښې له ټولو طبيعي عددونو
\emph{مخکښې} يو اضافي غړے ږدئ۔ لکه د هلبرټ د هوټل په استدلال کښې
چې مو په \olref[sfr][infinite][hilbert]{sec} کښې وکړل، په بداهتي
ډول دا اضافي لومړنے غړے بايد د ټول ترتيب په ډول کښې توپير رانه ولي۔
له بلې خوا، په $\omega \ordplus  1$ کښې له ټولو طبيعي عددونو
\emph{وروسته} يو اضافي غړے ږدئ۔ دا بيا بيخي بل ډول څيز دے!

\end{document}
